\unitopen{46}
\begin{candos}
\item rozpoznać funkcję partykuł (\de{doch, mal, eben, halt, wohl, bloß, nämlich}) i użyć kilku z nich;
\item rozumieć i używać dziesięciu popularnych idiomów oraz rozpoznawać dalsze z kontekstu;
\item rozbijać długie słowa złożone i tworzyć nowe zgodnie z regułami słowotwórstwa;
\item wybierać rejestr: potoczny, neutralny, formalny.
\end{candos}

\section{Słownictwo}
\vocabtable{46}

\section{Partykuły modalne}
Partykuły nie mają stałego tłumaczenia; zmieniają \emph{ton}. Najczęstsze:

\begin{center}\small
\begin{tabular}{@{}lp{0.38\linewidth}p{0.4\linewidth}@{}}
\toprule
\textsf{\textbf{Partykuła}} & \textsf{\textbf{Funkcja}} & \textsf{\textbf{Przykład}}\\
\midrule
\de{doch} & przypomina o czymś oczywistym; kontr. & \de{Das ist doch klar.} \de{Du kommst doch mit?}\\
\de{mal} & łagodzi prośbę & \de{Kannst du mal kurz helfen?}\\
\de{eben}, \de{halt} & „po prostu tak jest” & \de{Das ist eben so.} \de{Das ist halt so.}\\
\de{ja} & „jak wiadomo” & \de{Das wissen Sie ja.}\\
\de{wohl} & przypuszczenie & \de{Er ist wohl krank.}\\
\de{bloß} & napięcie, groźba & \de{Mach das bloß nicht!}\\
\de{nämlich} & wyjaśnienie & \de{Ich komme später, ich habe nämlich einen Termin.}\\
\bottomrule
\end{tabular}
\end{center}

\section{Idiomy}
\de{Ich drücke dir die Daumen.} (trzymam kciuki) \de{Das ist nicht mein Bier.} (to nie moja sprawa) \de{Ich verstehe nur Bahnhof.} (nic
nie rozumiem) \de{Das ist ein Kinderspiel.} (dziecinnie proste) \de{Wir sitzen im selben Boot.} (jesteśmy w tej samej sytuacji) \de{Das
Eis brechen} (przełamać lody) \de{Auf dem Holzweg sein} (iść złą drogą) \de{Das kannst du vergessen.} (zapomnij o tym).
Nie uczysz się idiomów hurtem: dziesięć najużyteczniejszych w Anki plus rozpoznawanie reszty w tekście.

\section{Słowotwórstwo}
Wyrazy złożone: \de{die Datenbank} $+$ \de{der Server} $=$ \de{der Datenbankserver}; rodzaj wyznacza ostatni człon. Złożenia z \de{-s-} (\de{Arbeitsvertrag}), z \de{-n-}
(\de{Kartenzahlung}), bez łącznika (\de{Tischlampe}). Przedrostki: \de{ver-} (\de{verbessern}), \de{ent-} (\de{entwickeln}), \de{er-} (\de{erledigen}),
\de{be-} (\de{bearbeiten}), \de{aus-}, \de{ein-}, \de{über-}. Przyrostki rzeczowników: \de{-ung}, \de{-heit}, \de{-keit}, \de{-schaft}, \de{-er} (wykonawca). Przymiotniki: \de{-bar}, \de{-lich}, \de{-ig}, \de{-los}.

\begin{zasada}[Czytanie długiego słowa]
Rozbij od końca: \de{Donaudampfschifffahrtsgesellschaft} to \de{Gesellschaft} (kto?) $\leftarrow$ \de{Schifffahrt} $\leftarrow$ \de{Dampf} $\leftarrow$ \de{Donau}.
Znaczenie bierzesz z końca, szczegóły z początku. W zawodowym tekście: \de{Datenschutzbeauftragter} $=$ \de{Daten} $+$ \de{Schutz} $+$ \de{Beauftragter}.
\end{zasada}

\begin{falle}
Polski nie ma partykuł modalnych, więc zdania brzmią niemiecko, ale płasko. Błąd: nadużywanie (po trzy w zdaniu) albo użycie \de{doch}
w znaczeniu angielskiego \emph{though}. Drugi: łączenie wyrazów złożonych spacją w stylu angielskim (\de{Daten Bank} zamiast \de{Datenbank}); to błąd
ortograficzny. Trzeci: używanie idiomów, których nie jesteś pewien. W razie wątpliwości mów prosto.
\end{falle}

\begin{szwajcaria}
Szwajcarski niemiecki ma własne partykuły i skrótowość: \de{gäll}, \de{oder}, \de{halt}, \de{eifach}. Szwajcarzy mówią \de{Das isch no guet}
(„to jest w porządku”). Rozumienie idiomów dialektalnych pomaga w integracji, ale nie wymuszaj mówienia dialektem.
\end{szwajcaria}

\begin{praca}
\de{Kannst du mal kurz auf meinen PR schauen?} \de{Das ist halt Legacy-Code.} \de{Das Ticket ist doch schon seit Wochen offen.} \de{Wir sitzen im selben
Boot, was den Termin betrifft.} Zwróć uwagę na rejestr: partykuły i idiomy są odpowiednie w zespole, nie w piśmie do klienta.
\end{praca}

\section{Dialogi}
\textbf{Dialog 1 --- rozmowa potoczna}
\begin{dialog}
\turn{A}{Kannst du mal kurz auf meinen Pull Request schauen? Ich verstehe nur Bahnhof, was die Fehlermeldung betrifft.}{Możesz rzucić okiem na mój pull request? Nic nie rozumiem z komunikatu o błędzie.}
\turn{B}{Klar. Das ist doch ein bekanntes Problem. Das liegt halt am Cache.}{Jasne. To przecież znany problem. Chodzi po prostu o cache.}
\turn{A}{Ach so, das hätte ich wohl wissen müssen.}{Ach tak, powinienem to chyba wiedzieć.}
\turn{B}{Kein Problem, ich drücke dir die Daumen, dass es jetzt klappt.}{Nie ma problemu, trzymam kciuki, żeby teraz zadziałało.}
\end{dialog}

\section{Ćwiczenia}
\begin{cwiczenia}
\exitem{Złagodź prośbę partykułą: \de{Kannst du mir helfen?}}{\de{Kannst du mir mal helfen?}}
\exitem{Przypomnij oczywistość: \de{Das ist klar.} (z partykułą)}{\de{Das ist doch klar.}}
\exitem{Wyraź przypuszczenie: \de{Er ist krank.} (z partykułą)}{\de{Er ist wohl krank.}}
\exitem{Wyjaśnij przyczynę: \de{Ich komme später. Ich habe einen Termin.} (nämlich)}{\de{Ich komme später, ich habe nämlich einen Termin.}}
\exitem{Idiom: „Trzymam kciuki.”}{\de{Ich drücke dir die Daumen.}}
\exitem{Idiom: „To nie moja sprawa.”}{\de{Das ist nicht mein Bier.}}
\exitem{Idiom: „Nic nie rozumiem.”}{\de{Ich verstehe nur Bahnhof.}}
\exitem{Rozbij wyraz: \de{Datenbankserver}.}{\de{Daten} + \de{Bank} + \de{Server}; rodzaj: \de{der Server} $\rightarrow$ \de{der Datenbankserver}}
\exitem{Popraw ortografię: \de{Daten Bank Server}.}{\de{Datenbankserver}}
\exitem{Który rejestr: \de{Das ist halt so} (formalny czy potoczny)?}{potoczny}
\end{cwiczenia}

\section{Trening słownictwa (generowany)}
\drill{46}

\weekend{Przeczytaj artykuł z SRF lub \de{Die Zeit} i wypisz dziesięć partykuł modalnych z kontekstem, pięć wyrazów złożonych (rozbij) i trzy idiomy.
Napisz dwa krótkie e-maile: potoczny do kolegi i formalny do klienta o tym samym.}
